\documentclass[11pt,a4paper]{article}

\usepackage[T1]{fontenc}
\usepackage[utf8]{inputenc}
\usepackage{lmodern}
\usepackage{microtype}
\usepackage[margin=2.8cm]{geometry}
\usepackage{amsmath,amssymb,amsthm,mathtools}
\usepackage{enumitem}
\usepackage{booktabs}
\usepackage{graphicx}
\usepackage{listings}
\usepackage{hyperref}

\hypersetup{
  colorlinks=true,
  linkcolor=blue,
  citecolor=blue,
  urlcolor=blue
}

\setlist[itemize]{leftmargin=2em}
\setlist[enumerate]{leftmargin=2em}
\lstset{
  basicstyle=\ttfamily\small,
  breaklines=true,
  columns=fullflexible,
  frame=single,
  language=Python,
  showstringspaces=false
}

\newtheorem{theorem}{Theorem}[section]
\newtheorem{proposition}[theorem]{Proposition}
\newtheorem{lemma}[theorem]{Lemma}
\newtheorem{definition}[theorem]{Definition}
\newtheorem{remark}[theorem]{Remark}

\newcommand{\R}{\mathbb{R}}
\newcommand{\N}{\mathbb{N}}
\newcommand{\E}{\mathbb{E}}
\newcommand{\Prob}{\mathbb{P}}
\newcommand{\dd}{\,\mathrm{d}}
\newcommand{\HS}{\mathrm{HS}}
\newcommand{\Id}{\mathrm{Id}}
\newcommand{\norm}[1]{\left\lVert #1 \right\rVert}
\newcommand{\ip}[2]{\left\langle #1,#2 \right\rangle}
\newcommand{\dom}{D}

\title{Numerical Approximation of a Stochastic Heat Equation}
\author{Seminar Notes}
\date{\today}

\begin{document}
\maketitle

\begin{abstract}
These notes introduce the one-dimensional stochastic heat equation with additive
space-time white noise and explain why its mathematically useful formulation is
an abstract stochastic evolution equation on the Hilbert space \(L^2(0,1)\).
After recalling the deterministic heat equation, semigroups, Wiener processes,
cylindrical noise and mild solutions, the notes derive the spectral
representation of the solution. The numerical part focuses on spectral Galerkin
truncation and exact simulation of the resulting Ornstein--Uhlenbeck modes. The
main convergence result proves mean-square convergence of the spectral
approximation and gives the rate \(N^{-1/2}\) for the root mean-square error of
the stochastic truncation part.
\end{abstract}
\paragraph*{Note on the use of AI tools.}
During the preparation of these notes, AI-assisted tools were used in a
supporting role. ChatGPT was used for the linguistic revision of individual
formulations and for translation into English. The code used to generate the
numerical plots was created with the help of Codex and was subsequently
reviewed, adapted, and validated.

\tableofcontents

\section{Introduction}

The deterministic heat equation is a basic model for diffusion. On a spatial
domain \(\Omega\), the unknown \(u(t,x)\) describes a temperature, concentration,
or density. In one space dimension the sign of the second derivative gives a
direct intuition for the equation. At a point where the graph is locally
concave, \(u_{xx}<0\), so the equation \(u_t=u_{xx}\) makes the value decrease
in time. At a point where the graph is locally convex, \(u_{xx}>0\), so the
value increases. Peaks are therefore lowered and valleys are filled, which is
the basic smoothing effect of heat flow. A derivation from energy balance and
Fourier's law is given in the PDE I notes \cite{stinnerPDE1}. In the simplest
homogeneous case one obtains
\[
  \partial_t u(t,x)=\Delta u(t,x).
\]
The role of initial and boundary conditions is also part of this deterministic
theory \cite{stinnerPDE1}.

In many applications the system is not perfectly deterministic. Random
microscopic fluctuations, unresolved forcing, or measurement errors can be
modeled by adding noise to the equation. The formal model studied here is
\[
  \partial_t u(t,x)=\partial_{xx}u(t,x)+\dot W(t,x),
  \qquad x\in(0,1),\quad t>0,
\]
with homogeneous Dirichlet boundary conditions
\[
  u(t,0)=u(t,1)=0.
\]
The symbol \(\dot W\) denotes space-time white noise. It is not a classical
function and, even as a random distribution, it is not an \(L^2(0,1)\)-valued
forcing term. Therefore the equation cannot be understood pointwise in the same
sense as a smooth deterministic PDE. One could try to interpret the equation
pathwise for fixed \(\omega\) in a distributional sense. In these notes we
instead deliberately use the stochastic semigroup formulation, because it is a
robust framework for more general SPDEs and is very convenient for numerical
analysis: the same formula that defines the solution also leads directly to a
convergent spectral algorithm \cite{delSarto2025,daPratoZabczyk,lordPowellShardlow}.
The goal of these notes is to show how the equation is formulated and solved in
the Hilbert space
\[
  H=L^2(0,1),
\]
and how this formulation directly leads to a practical and convergent numerical
method.

The main information is taken from lecture notes for courses held at TU
Darmstadt and EPFL. The deterministic heat equation and eigenfunction material
are based on the PDE I notes \cite{stinnerPDE1}; the semigroup and abstract
Cauchy-problem background is based on the PDE II notes \cite{stinnerPDE2}; the
stochastic analysis, Wiener processes in Hilbert spaces, stochastic convolution
and additive SPDE framework are based on Del Sarto's lecture notes
\cite{delSarto2025}. The coordinate-free viewpoint on white noise and
stochastic PDEs is taken as supplementary background from Hairer's EPFL notes
\cite{hairer2009}. Further standard references used for SPDE theory and
numerical SPDEs are \cite{daPratoZabczyk,lordPowellShardlow}.

\section{The Deterministic Heat Equation}

\subsection{Classical formulation}

On a smooth domain \(\Omega\subset\R^d\), the inhomogeneous heat equation has the
form
\[
  \partial_t u(t,x)=\Delta u(t,x)+f(t,x),
  \qquad (t,x)\in(0,T]\times\Omega.
\]
To obtain a well-posed problem one prescribes an initial condition
\[
  u(0,x)=u_0(x)
\]
and suitable boundary conditions. In this text we use homogeneous Dirichlet
conditions on \((0,1)\):
\[
  u(t,0)=u(t,1)=0.
\]
This is the natural deterministic prototype for the stochastic heat equation
treated later.

\subsection{Heat flow on the real line}

On \(\R^d\), the fundamental solution of the heat equation is
\[
  \Phi(x,t)=\frac{1}{(4\pi t)^{d/2}}\exp\left(-\frac{|x|^2}{4t}\right),
  \qquad t>0.
\]
For bounded and continuous initial data \(u_0\), the solution of the Cauchy
problem on \(\R^d\) is given by convolution,
\[
  u(t,x)=\int_{\R^d}\Phi(x-y,t)u_0(y)\dd y.
\]
This formula illustrates the smoothing effect of the heat equation. Even if the
initial data is only moderately regular, the solution becomes smooth for
positive times. This smoothing property is also the reason why the stochastic
convolution below is well-defined in one spatial dimension. The fundamental
solution formula and its smoothing interpretation are standard material from
the deterministic heat equation theory in \cite{stinnerPDE1}.

\subsection{\texorpdfstring{Dirichlet Laplacian on \((0,1)\)}{Dirichlet Laplacian on (0,1)}}

For the interval \((0,1)\), let
\[
  H=L^2(0,1),\qquad
  \ip{u}{v}_H=\int_0^1 u(x)v(x)\dd x.
\]
The Dirichlet Laplace operator is
\[
  A=\partial_{xx},\qquad
  \dom(A)=H^2(0,1)\cap H_0^1(0,1).
\]
It is negative self-adjoint and has compact resolvent. In particular, it admits
an orthonormal eigenbasis in \(H\). Explicitly,
\[
  e_k(x)=\sqrt{2}\sin(k\pi x),
  \qquad
  \lambda_k=(k\pi)^2,
  \qquad k\in\N,
\]
and
\[
  Ae_k=-\lambda_k e_k.
\]
Thus every \(v\in H\) has the Fourier expansion
\[
  v=\sum_{k=1}^{\infty} v_k e_k,
  \qquad
  v_k=\ip{v}{e_k}_H,
\]
with
\[
  \norm{v}_H^2=\sum_{k=1}^{\infty}|v_k|^2.
\]
The general existence of eigenvalues and an orthonormal eigenbasis for the
Dirichlet Laplacian on bounded domains is part of the variational theory
presented in \cite{stinnerPDE1}. The explicit sine basis is the special
one-dimensional formula on the interval.

\section{Semigroups and Abstract Formulation}

\subsection{\texorpdfstring{\(C_0\)-semigroups}{C0-semigroups}}

The deterministic heat equation can be written as an ordinary differential
equation on an infinite-dimensional function space. This requires semigroup
theory.

\begin{definition}
Let \(X\) be a Banach space. A family \((S(t))_{t\ge0}\subset\mathcal L(X)\) is
called a semigroup if
\[
  S(0)=\Id,
  \qquad
  S(t)S(s)=S(t+s)
  \quad\text{for all }s,t\ge0.
\]
It is called strongly continuous, or a \(C_0\)-semigroup, if
\[
  \lim_{t\downarrow0}S(t)x=x
  \qquad\text{for all }x\in X.
\]
\end{definition}

If \(A\) is the generator of a \(C_0\)-semigroup, then the abstract Cauchy problem
\[
  \frac{\dd}{\dd t}u(t)=Au(t),
  \qquad u(0)=u_0,
\]
has the solution
\[
  u(t)=S(t)u_0.
\]
This is the infinite-dimensional analogue of \(u(t)=e^{tA}u_0\) for a matrix
\(A\). The definitions and generator results used here are developed in the PDE
II notes \cite{stinnerPDE2}.

\subsection{The heat semigroup}

For the Dirichlet Laplacian \(A=\partial_{xx}\) on \(H=L^2(0,1)\), define
\[
  S(t)=e^{tA}.
\]
Using the eigenbasis \((e_k)_{k\in\N}\), the semigroup acts by
\[
  S(t)v
  =
  \sum_{k=1}^{\infty}e^{-\lambda_k t}v_k e_k.
\]
In particular,
\[
  S(t)e_k=e^{-\lambda_k t}e_k.
\]
For every \(t\ge0\), \(S(t)\) is a contraction on \(H\):
\[
  \norm{S(t)v}_H^2
  =
  \sum_{k=1}^{\infty}e^{-2\lambda_k t}|v_k|^2
  \le
  \norm{v}_H^2.
\]
The deterministic homogeneous heat equation with Dirichlet boundary conditions
can therefore be written as
\[
  u'(t)=Au(t),
  \qquad u(0)=u_0,
\]
with solution
\[
  u(t)=S(t)u_0.
\]

\subsection{Inhomogeneous deterministic equation}

Before turning to noise it is useful to recall the deterministic inhomogeneous
problem
\[
  u'(t)=Au(t)+f(t),
  \qquad u(0)=u_0.
\]
For a deterministic source term \(f\), the variation-of-constants formula gives
\[
  u(t)=S(t)u_0+\int_0^t S(t-s)f(s)\dd s.
\]
In PDE notation this corresponds, for example, to
\[
  \partial_t u(t,x)=\partial_{xx}u(t,x)+f(t,x)
\]
with the same initial and boundary conditions as above. Thus the mild
formulation is not a new idea introduced only for stochastic equations; in the
deterministic setting it is precisely the variation-of-constants formula from
semigroup theory \cite{stinnerPDE2}. In the stochastic setting the same formula
is used as the guiding analogy, but the deterministic integral is replaced by a
stochastic integral with respect to a Wiener process. Giving meaning to this
stochastic integral is the next step.

\section{Brownian Motion, Wiener Processes and Cylindrical Noise}

\subsection{Brownian motion}

\begin{definition}
A real-valued stochastic process \((\beta(t))_{t\ge0}\) is a standard Brownian
motion if
\begin{enumerate}
  \item \(\beta(0)=0\) almost surely,
  \item \(\beta(t)-\beta(s)\sim\mathcal N(0,t-s)\) for \(0\le s<t\),
  \item the increments are independent,
  \item the paths are continuous almost surely.
\end{enumerate}
\end{definition}

Brownian motion is also called a Wiener process. It is the basic driving noise
for finite-dimensional stochastic differential equations; see
\cite{delSarto2025}.

\subsection{Hilbert-space setup}

We work on
\[
  H=L^2(0,1)
\]
with orthonormal basis
\[
  e_k(x)=\sqrt{2}\sin(k\pi x),
  \qquad k\in\N.
\]
Let \((\beta_k)_{k\in\N}\) be independent real standard Brownian motions on a
filtered probability space
\[
  (\Omega,\mathcal F,(\mathcal F_t)_{t\ge0},\Prob).
\]

\subsection{Cylindrical Wiener process}

One would like to write
\[
  W(t)=\sum_{k=1}^{\infty}\beta_k(t)e_k.
\]
However, this series does not converge in \(H\). Indeed, by orthonormality,
\[
  \E\norm{W(t)}_H^2
  =
  \sum_{k=1}^{\infty}\E|\beta_k(t)|^2
  =
  \sum_{k=1}^{\infty} t
  =
  \infty.
\]
Thus \(W(t)\) is not an \(H\)-valued random variable. This is the main
mathematical reason why space-time white noise is singular.

\begin{definition}
Let \((e_k)_{k\in\N}\) be an orthonormal basis of \(H\) and let
\((\beta_k)_{k\in\N}\) be independent real Brownian motions. The formal process
\[
  W(t)=\sum_{k=1}^{\infty}\beta_k(t)e_k
\]
is called a cylindrical Wiener process on \(H\). Equivalently, it is defined by
its action on test functions \(h\in H\):
\[
  W_t(h)
  =
  \sum_{k=1}^{\infty}\ip{h}{e_k}_H\beta_k(t).
\]
\end{definition}

For fixed \(h\in H\), this series converges in \(L^2(\Omega)\), since
\[
  \E|W_t(h)|^2
  =
  t\sum_{k=1}^{\infty}|\ip{h}{e_k}_H|^2
  =
  t\norm{h}_H^2.
\]
Hence
\[
  W_t(h)\sim\mathcal N(0,t\norm{h}_H^2).
\]
Moreover, for \(g,h\in H\),
\[
\begin{aligned}
  \E\bigl[W_t(h)W_s(g)\bigr]
  &=
  \E\left[
    \sum_{k=1}^{\infty}\ip{h}{e_k}_H\beta_k(t)
    \sum_{\ell=1}^{\infty}\ip{g}{e_\ell}_H\beta_\ell(s)
  \right]  \\
  &=
  \sum_{k,\ell=1}^{\infty}
  \ip{h}{e_k}_H\ip{g}{e_\ell}_H
  \E[\beta_k(t)\beta_\ell(s)]  \\
  &=
  (t\wedge s)
  \sum_{k=1}^{\infty}\ip{h}{e_k}_H\ip{g}{e_k}_H  \\
  &=
  (t\wedge s)\ip{h}{g}_H,
\end{aligned}
\]
where independence gives
\(\E[\beta_k(t)\beta_\ell(s)]=(t\wedge s)\delta_{k\ell}\), and the last
identity is Parseval's formula.
This is the rigorous way to treat the formal series when the covariance
operator is the identity. In the terminology of \cite{delSarto2025}, it is a
generalized Wiener process with covariance \(Q=\Id\), hence a cylindrical Wiener
process.

\begin{remark}
Hairer's notes define this object in a coordinate-free way: \(W_t\) is a
Gaussian linear map on test functions such that
\[
  \E[W_t(h)W_s(g)]=(t\wedge s)\ip{h}{g}_H
\]
for all \(g,h\in H\) \cite{hairer2009}. The construction in
\cite{delSarto2025} uses an orthonormal basis to realize the same object by
the series \(\sum_k\beta_k(t)e_k\). The two viewpoints are equivalent for the
present purpose, because the law of \(W_t(h)\) depends only on
\(\norm{h}_H\) and the covariance formula above, not on the particular
orthonormal basis used to represent \(h\).
\end{remark}

\section{The Stochastic Heat Equation as an Abstract SPDE}

\subsection{Abstract formulation}

As in the deterministic case, the starting point is an abstract Cauchy problem.
The difference is that the forcing is now random. Formally, the stochastic heat
equation becomes a stochastic differential equation on the Hilbert space
\(H=L^2(0,1)\). This notation should still be read with care: the time
derivative of the driving noise is only a distributional object, and \(W(t)\)
itself is cylindrical rather than \(H\)-valued \cite{delSarto2025,hairer2009}.

Let
\[
  H=L^2(0,1),
  \qquad
  A=\partial_{xx},
  \qquad
  \dom(A)=H^2(0,1)\cap H_0^1(0,1).
\]
Let \(W\) be a cylindrical Wiener process on \(H\). The stochastic heat equation
with additive space-time white noise is formally
\[
  \dd u(t)=Au(t)\dd t+\dd W_t,
  \qquad u(0)=u_0.
\]
In PDE notation this corresponds to
\[
  \partial_t u(t,x)=\partial_{xx}u(t,x)+\dot W(t,x),
  \qquad x\in(0,1),\quad t>0,
\]
with boundary conditions
\[
  u(t,0)=u(t,1)=0.
\]

\subsection{Why the classical formulation is not sufficient}

There are two related obstacles. First, the cylindrical Wiener process \(W(t)\)
is not an \(H\)-valued random variable. Second, the formal derivative
\(\dot W(t,x)\) is not a classical function but only a distributional symbol.
Consequently, the expression
\[
  Au(t)+\dot W(t)
\]
is not meaningful as an \(H\)-valued classical right-hand side. One therefore
does not first search for a pointwise differentiable solution. The appropriate
object is a mild solution, obtained by integrating the equation against the heat
semigroup \cite{delSarto2025,daPratoZabczyk}.

\section{Mild Solutions}

Let \(S(t)=e^{tA}\) be the heat semigroup generated by \(A\).

\subsection{Stochastic integral with cylindrical noise}

Because \(W\) is cylindrical rather than \(H\)-valued, not every bounded
operator can be integrated against \(W\). The smoothing of the heat semigroup
enters through the Hilbert-Schmidt norm. If \(B\in\mathcal L(H)\), then
\[
  \norm{B}_{\HS}^2
  =
  \sum_{k=1}^{\infty}\norm{Be_k}_H^2
\]
whenever this sum is finite. The value is independent of the chosen
orthonormal basis. Such operators are called Hilbert-Schmidt operators; this is
the operator class used to integrate against cylindrical Wiener processes
\cite{delSarto2025,daPratoZabczyk}.

For a deterministic operator-valued function
\(\Phi:[0,t]\to\mathcal L(H)\) satisfying
\[
  \int_0^t\norm{\Phi(s)}_{\HS}^2\dd s<\infty,
\]
the stochastic integral is defined by the \(L^2(\Omega;H)\)-limit
\[
  \int_0^t \Phi(s)\dd W_s
  =
  \lim_{M\to\infty}
  \sum_{k=1}^{M}\int_0^t \Phi(s)e_k\dd\beta_k(s).
\]
The Ito isometry for cylindrical noise, stated in Appendix
\ref{app:ito-isometry}, gives
\[
  \E\left\|\int_0^t \Phi(s)\dd W_s\right\|_H^2
  =
  \int_0^t\norm{\Phi(s)}_{\HS}^2\dd s.
\]

\begin{remark}
For random integrands one also requires the usual measurability and
adaptedness, or predictability, assumptions so that the integrand only uses
information available up to the current time. In the linear additive equation
treated here the integrand \(S(t-s)\) is deterministic, so this issue is
harmless. It becomes essential for multiplicative noise.
\end{remark}

\begin{remark}
Hairer's notes often formulate stochastic integrals and white noise without
choosing a distinguished orthonormal basis, starting instead from Gaussian
families indexed by test functions \cite{hairer2009}. The basis construction
above is the concrete Hilbert-space version used in \cite{delSarto2025}. The
definitions agree because the resulting integrals have the same covariance
structure, as expressed by the Ito isometry.
\end{remark}

\subsection{Definition}

\begin{definition}
An \(H\)-valued adapted process \((u(t))_{t\ge0}\) is called a mild solution of
the stochastic heat equation if, for every \(t\ge0\),
\[
  u(t)
  =
  S(t)u_0
  +
  \int_0^t S(t-s)\dd W_s.
\]
The stochastic integral
\[
  \int_0^t S(t-s)\dd W_s
\]
is called the stochastic convolution.
\end{definition}

This formula is the stochastic analogue of the deterministic
variation-of-constants formula. The stochastic convolution is the central object
in the theory of linear SPDEs with additive noise; see \cite{delSarto2025}.

\begin{remark}
For the linear additive equation, the weak and mild formulations are
equivalent under the standard assumptions on \(A\), \(S(t)\), \(u_0\), and the
stochastic convolution \cite{daPratoZabczyk,delSarto2025}. The mild
formulation is the natural one for this numerical seminar because it takes the
variation-of-constants formula as the definition of the solution and leads
immediately to the spectral Galerkin method below. Weak formulations are often
more intrinsic analytically, but they are less direct for implementing the
modal approximation.
\end{remark}

\section{Spectral Decomposition and Existence of the Solution}

\subsection{Fourier representation}

Write
\[
  u_0=\sum_{k=1}^{\infty}u_{0,k}e_k,
  \qquad
  u_{0,k}=\ip{u_0}{e_k}_H.
\]
The deterministic part is
\[
  S(t)u_0
  =
  \sum_{k=1}^{\infty}e^{-\lambda_k t}u_{0,k}e_k.
\]
For the stochastic convolution, the definition of the integral with respect to
the cylindrical Wiener process gives
\[
\begin{aligned}
  \int_0^t S(t-s)\dd W_s
  &=
  \lim_{M\to\infty}
  \sum_{k=1}^{M}
  \int_0^t S(t-s)e_k\dd\beta_k(s)  \\
  &=
  \lim_{M\to\infty}
  \sum_{k=1}^{M}
  \left(\int_0^t e^{-\lambda_k(t-s)}\dd\beta_k(s)\right)e_k,
\end{aligned}
\]
where the convergence is in \(L^2(\Omega;H)\). Hence the mild solution has the
spectral representation
\[
  u(t)
  =
  \sum_{k=1}^{\infty}
  \left(
    e^{-\lambda_k t}u_{0,k}
    +
    \int_0^t e^{-\lambda_k(t-s)}\dd\beta_k(s)
  \right)e_k.
\]
Equivalently, if \(u(t)=\sum_k u_k(t)e_k\), then comparison of the Fourier
coefficients gives
\[
  u_k(t)
  =
  e^{-\lambda_k t}u_{0,k}
  +
  \int_0^t e^{-\lambda_k(t-s)}\dd\beta_k(s).
\]
This is exactly the mild form of the scalar Ornstein--Uhlenbeck equation
\[
  \dd u_k(t)=-\lambda_k u_k(t)\dd t+\dd\beta_k(t),
  \qquad u_k(0)=u_{0,k}.
\]
The mode-wise interpretation is standard for linear additive SPDEs diagonalized
by the generator eigenbasis \cite{delSarto2025,lordPowellShardlow}.

\subsection{Existence theorem}

\begin{theorem}[Existence of the mild solution in \(L^2(0,1)\)]
Let \(H=L^2(0,1)\), let \(A=\partial_{xx}\) with homogeneous Dirichlet boundary
conditions, and let \(W\) be a cylindrical Wiener process on \(H\). For every
\(u_0\in H\) and every \(t>0\),
\[
  u(t)=S(t)u_0+\int_0^t S(t-s)\dd W_s
\]
is a well-defined \(H\)-valued random variable. Moreover,
\[
  \E\norm{u(t)}_H^2<\infty.
\]
\end{theorem}

\begin{proof}
Step 1: well-defined deterministic and stochastic terms. The deterministic
term is well-defined since \(S(t)\) is a contraction on \(H\):
\[
  \norm{S(t)u_0}_H\le \norm{u_0}_H.
\]
For the stochastic term we check the Hilbert-Schmidt condition. Since
\(S(r)e_k=e^{-\lambda_k r}e_k\),
\[
  \norm{S(r)}_{\HS}^2
  =
  \sum_{k=1}^{\infty}\norm{S(r)e_k}_H^2
  =
  \sum_{k=1}^{\infty}e^{-2\lambda_k r}.
\]
Therefore
\[
  \int_0^t\norm{S(r)}_{\HS}^2\dd r
  =
  \sum_{k=1}^{\infty}\int_0^t e^{-2\lambda_k r}\dd r
  =
  \sum_{k=1}^{\infty}\frac{1-e^{-2\lambda_k t}}{2\lambda_k}.
\]
Since \(\lambda_k=(k\pi)^2\),
\[
  \sum_{k=1}^{\infty}\frac{1}{\lambda_k}
  =
  \frac{1}{\pi^2}\sum_{k=1}^{\infty}\frac{1}{k^2}
  <\infty.
\]
Thus \(\int_0^t S(t-s)\dd W_s\) is well-defined as an \(H\)-valued stochastic
integral.

Step 2: second-moment estimate. Define
\[
  Z(t)=\int_0^t S(t-s)\dd W_s
  =
  \sum_{k=1}^{\infty}
  \left(
    \int_0^t e^{-\lambda_k(t-s)}\dd\beta_k(s)
  \right)e_k.
\]
By the Ito isometry,
\[
  \E\left|
    \int_0^t e^{-\lambda_k(t-s)}\dd\beta_k(s)
  \right|^2
  =
  \int_0^t e^{-2\lambda_k(t-s)}\dd s
  =
  \frac{1-e^{-2\lambda_k t}}{2\lambda_k}.
\]
Therefore
\[
  \E\norm{Z(t)}_H^2
  =
  \sum_{k=1}^{\infty}
  \frac{1-e^{-2\lambda_k t}}{2\lambda_k}
  \le
  \frac12\sum_{k=1}^{\infty}\frac{1}{\lambda_k}
  <\infty.
\]
The convergence of the series follows from \(\lambda_k=(k\pi)^2\).

Step 3: combination and uniqueness. Since the stochastic convolution is
centered, the deterministic and stochastic parts give
\[
  \E\norm{u(t)}_H^2
  =
  \norm{S(t)u_0}_H^2
  +
  \E\norm{Z(t)}_H^2
  <\infty.
\]
The mild solution is unique in this class, because the formula determines all
Fourier coefficients uniquely; equivalently, each coefficient is the unique
solution of its scalar Ornstein--Uhlenbeck equation. This is the standard
linear additive SPDE uniqueness argument \cite{delSarto2025,daPratoZabczyk}.
\end{proof}

\begin{remark}[Spatial regularity]
The same computation also shows the limited spatial regularity caused by white
noise. For \(u_0=0\) and \(\gamma\ge0\),
\[
  \E\norm{(-A)^\gamma Z(t)}_H^2
  =
  \sum_{k=1}^{\infty}
  \lambda_k^{2\gamma}
  \frac{1-e^{-2\lambda_k t}}{2\lambda_k}.
\]
Since \(\lambda_k=(k\pi)^2\), the summand behaves like
\(k^{4\gamma-2}\). The series is finite precisely for
\(\gamma<1/4\). Thus the stochastic part is in
\(D((-A)^\gamma)\) for every \(\gamma<1/4\), but in general not for
\(\gamma\ge1/4\). In Sobolev language this corresponds to spatial regularity
strictly below \(H^{1/2}(0,1)\). This regularity threshold is a standard
feature of the one-dimensional stochastic heat equation with additive white
noise \cite{daPratoZabczyk,hairer2009}.
\end{remark}

\begin{remark}[Higher spatial dimensions]
On a bounded domain in dimension \(d\), the proof changes at exactly the point
where the eigenvalue sum is estimated. By Weyl's law,
\(\lambda_j\sim Cj^{2/d}\). The condition for an \(L^2\)-valued stochastic
convolution with space-time white noise becomes
\[
  \sum_{j=1}^{\infty}\frac1{\lambda_j}<\infty,
\]
which holds only for \(d=1\). In \(d=2\) the corresponding series behaves like
\(\sum_j j^{-1}\) and diverges logarithmically; for \(d\ge3\) it diverges even
more strongly. Thus the one-dimensional proof does not extend to an
\(L^2\)-valued solution driven by spatially white noise in higher dimensions.
One then has to use spatially colored noise, weaker distribution-valued
solution spaces, or, for singular nonlinear equations, renormalization
\cite{daPratoZabczyk,hairer2009}.
\end{remark}

\section{Numerical Simulation and Convergence}

\subsection{Spectral Galerkin approximation}

The numerical idea is to approximate the infinite Fourier expansion by keeping
only the first \(N\) eigenfunctions. This replaces the infinite-dimensional
SPDE by a finite-dimensional system of \(N\) stochastic differential equations.
In the sine basis these equations are decoupled, because the basis diagonalizes
the operator \(A\). For \(N\in\N\), the Galerkin space is
\[
  H_N=\operatorname{span}\{e_1,\dots,e_N\}
\]
and let \(P_N:H\to H_N\) be the orthogonal projection
\[
  P_Nv=\sum_{k=1}^{N}\ip{v}{e_k}_H e_k.
\]
The spectral Galerkin approximation is
\[
  \dd u^N(t)=A_Nu^N(t)\dd t+P_N\dd W_t,
  \qquad
  u^N(0)=P_Nu_0,
\]
where \(A_N=P_NA\). Its mild solution is
\[
  u^N(t)
  =
  S(t)P_Nu_0
  +
  \int_0^t S(t-s)P_N\dd W_s.
\]
Equivalently,
\[
  u^N(t,x)=\sum_{k=1}^{N}u_k(t)e_k(x),
\]
where
\[
  \dd u_k(t)=-\lambda_k u_k(t)\dd t+\dd\beta_k(t),
  \qquad k=1,\dots,N.
\]

\subsection{Convergence theorem for spectral truncation}

\begin{theorem}[Mean-square convergence of spectral Galerkin approximation]
Let \(u(t)\) be the mild solution of the stochastic heat equation on \((0,1)\),
and let \(u^N(t)\) be its spectral Galerkin approximation. Then for every
\(t>0\),
\[
  \E\norm{u(t)-u^N(t)}_H^2\longrightarrow0
  \qquad\text{as }N\to\infty.
\]
If additionally \(u_0\in\dom((-A)^\alpha)\) for some \(\alpha>0\), then
\[
  \E\norm{u(t)-u^N(t)}_H^2
  \le
  C\lambda_{N+1}^{-2\alpha}\norm{(-A)^\alpha u_0}_H^2
  +
  C\sum_{k=N+1}^{\infty}\frac1{\lambda_k}.
\]
In particular, since \(\lambda_k=(k\pi)^2\),
\[
  \sum_{k=N+1}^{\infty}\frac1{\lambda_k}
  \le
  \frac{1}{\pi^2N},
\]
and \(\lambda_{N+1}^{-2\alpha}\le C N^{-4\alpha}\). Hence
\[
  \E\norm{u(t)-u^N(t)}_H^2
  \le
  C N^{-4\alpha}\norm{(-A)^\alpha u_0}_H^2
  +
  C N^{-1}.
\]
For \(u_0=0\), or if only the stochastic truncation part is considered, the
corresponding \(L^2(\Omega;H)\)-error is of order \(N^{-1/2}\).
\end{theorem}

\begin{proof}
Since \(P_N\) commutes with \(S(t)\) in the eigenbasis, we write
\[
  u(t)-u^N(t)
  =
  S(t)(I-P_N)u_0
  +
  \int_0^t S(t-s)(I-P_N)\dd W_s.
\]
We now estimate the deterministic and stochastic parts separately.

\emph{Step 1: convergence of the deterministic part.}
Write
\[
  u_0=\sum_{k=1}^{\infty}u_{0,k}e_k.
\]
Then
\[
  S(t)(I-P_N)u_0
  =
  \sum_{k=N+1}^{\infty}e^{-\lambda_k t}u_{0,k}e_k,
\]
and therefore
\[
  \norm{S(t)(I-P_N)u_0}_H^2
  =
  \sum_{k=N+1}^{\infty}e^{-2\lambda_k t}|u_{0,k}|^2
  \longrightarrow0.
\]
The convergence follows because the right-hand side is the tail of a
convergent nonnegative series dominated by
\(\sum_{k=1}^{\infty}|u_{0,k}|^2=\norm{u_0}_H^2\).

\emph{Step 2: rate for the deterministic part.}
If \(u_0\in\dom((-A)^\alpha)\), then by the \(\ell^2\)-norm of the Fourier
coefficients of \((-A)^\alpha u_0\),
\[
  \sum_{k=1}^{\infty}\lambda_k^{2\alpha}|u_{0,k}|^2<\infty.
\]
For \(k>N\), the eigenvalues are increasing and therefore
\(\lambda_k\ge\lambda_{N+1}\). Hence
\[
  1\le \lambda_{N+1}^{-2\alpha}\lambda_k^{2\alpha},
\]
which implies
\[
  |u_{0,k}|^2
  \le
  \lambda_{N+1}^{-2\alpha}\lambda_k^{2\alpha}|u_{0,k}|^2.
\]
Thus
\[
  \norm{S(t)(I-P_N)u_0}_H^2
  \le
  \lambda_{N+1}^{-2\alpha}
  \norm{(-A)^\alpha u_0}_H^2.
\]

\emph{Step 3: convergence and rate for the stochastic part.}
By the Ito isometry for cylindrical noise from Appendix
\ref{app:ito-isometry},
\[
  \E\left\|
    \int_0^t S(t-s)(I-P_N)\dd W_s
  \right\|_H^2
  =
  \int_0^t
  \norm{S(t-s)(I-P_N)}_{\HS}^2\dd s.
\]
For \(r>0\),
\[
  S(r)(I-P_N)e_k
  =
  \begin{cases}
    0, & k\le N,\\
    e^{-\lambda_k r}e_k, & k>N.
  \end{cases}
\]
Therefore
\[
  \norm{S(r)(I-P_N)}_{\HS}^2
  =
  \sum_{k=N+1}^{\infty}e^{-2\lambda_k r}.
\]
It follows that
\[
  \int_0^t
  \sum_{k=N+1}^{\infty}e^{-2\lambda_k r}\dd r
  =
  \sum_{k=N+1}^{\infty}
  \frac{1-e^{-2\lambda_k t}}{2\lambda_k}
  \le
  \frac12\sum_{k=N+1}^{\infty}\frac1{\lambda_k}.
\]
Since
\[
  \sum_{k=N+1}^{\infty}\frac1{\lambda_k}
  =
  \frac1{\pi^2}\sum_{k=N+1}^{\infty}\frac1{k^2}
  \le
  \frac1{\pi^2}\int_N^\infty \frac{\dd x}{x^2}
  =
  \frac1{\pi^2N},
\]
the stochastic truncation error tends to zero and satisfies the bound
\[
  \E\left\|
    \int_0^t S(t-s)(I-P_N)\dd W_s
  \right\|_H^2
  \le \frac{1}{2\pi^2N}.
\]

\emph{Step 4: combination of the errors.}
Finally,
\[
  \norm{a+b}_H^2\le2\norm{a}_H^2+2\norm{b}_H^2
\]
gives
\[
  \E\norm{u(t)-u^N(t)}_H^2
  \le
  2\norm{S(t)(I-P_N)u_0}_H^2
  +
  2\E\left\|
    \int_0^t S(t-s)(I-P_N)\dd W_s
  \right\|_H^2.
\]
Both terms converge to zero. The quantitative estimate follows from the two
bounds above.
\end{proof}

\begin{remark}[Higher-dimensional convergence]
As explained in the previous remark on higher spatial dimensions, the
eigenvalue sum \(\sum_j\lambda_j^{-1}\) diverges for spatially white noise when
\(d\ge2\). Therefore the stochastic convolution is not \(L^2\)-valued in that
case, and the above \(L^2(\Omega;H)\)-convergence proof does not carry over
directly. One must first change the setting, for example by using colored
noise or weaker solution spaces \cite{daPratoZabczyk,lordPowellShardlow}.
\end{remark}

\subsection{Time discretization}

For each mode \(k\le N\), the Galerkin system contains the scalar SDE
\[
  \dd u_k(t)=-\lambda_k u_k(t)\dd t+\dd\beta_k(t).
\]
Starting from this finite-dimensional SDE system, one would normally continue
with a time-stepping method such as Euler--Maruyama,
\[
  u_k^{n+1}
  =
  u_k^n-\lambda_k u_k^n\Delta t+\Delta\beta_k^n,
  \qquad
  \Delta\beta_k^n\sim\mathcal N(0,\Delta t).
\]
In the present linear additive problem, however, the exact solution of each
mode is known. Therefore the simpler and more accurate approach is to use the
exact Ornstein--Uhlenbeck transition
\[
  u_k(t_{n+1})
  =
  e^{-\lambda_k\Delta t}u_k(t_n)+\eta_k^n,
\]
where
\[
  \eta_k^n
  \sim
  \mathcal N\left(
    0,
    \frac{1-e^{-2\lambda_k\Delta t}}{2\lambda_k}
  \right).
\]
The random variables \(\eta_k^n\) are independent for different \(k\) and \(n\).
Thus the simulation is exact in time for the finite-dimensional Galerkin SPDE.

\subsection{Convergence with exact time stepping}

If the exact Ornstein--Uhlenbeck transition is used, then on the grid
\[
  t_n=n\Delta t
\]
there is no time discretization error for the finite-dimensional Galerkin
system. The only mathematical approximation error relative to the full SPDE is
the spectral truncation error:
\[
  u(t_n)-u^N(t_n).
\]
Consequently, the convergence theorem above fully describes the discretization
error of the implemented method at grid times.

\subsection{Implementation outline}

For fixed \(N\), final time \(T\), time step \(\Delta t\), and grid
\(t_n=n\Delta t\), the simulation can be implemented as follows.

\begin{enumerate}
  \item Compute \(\lambda_k=(k\pi)^2\) for \(k=1,\dots,N\).
  \item Compute the initial Fourier coefficients
  \[
    u_{0,k}=\int_0^1 u_0(x)e_k(x)\dd x.
  \]
  \item For every time step and every mode, draw
  \[
    \eta_k^n\sim
    \mathcal N\left(0,\frac{1-e^{-2\lambda_k\Delta t}}{2\lambda_k}\right).
  \]
  \item Update
  \[
    u_k^{n+1}=e^{-\lambda_k\Delta t}u_k^n+\eta_k^n.
  \]
  \item Reconstruct the physical-space approximation by
  \[
    u^N(t_n,x)=\sum_{k=1}^{N}u_k^n e_k(x).
  \]
\end{enumerate}

\section{Numerical Results}

This section reports a reproducible numerical experiment computed with the
spectral Ornstein--Uhlenbeck simulation described above. The path generation and
error computation code used for the table is included in Appendix
\ref{app:python-code}. The numerical method follows the exact modal transition
from the previous subsection and the standard spectral Galerkin viewpoint from
\cite{lordPowellShardlow}.

\subsection{Reference experiment}

A simple test case is
\[
  u_0=0,\qquad T=1.
\]
For each \(N\), simulate \(M\) independent paths of \(u^N(T)\). To estimate the
spectral error, use a larger reference dimension \(N_{\mathrm{ref}}\gg N\) and
drive both approximations with the same Brownian modes for \(k\le N\). Then
\[
  \left(
    \frac1M\sum_{m=1}^{M}
    \norm{u^{N_{\mathrm{ref}},(m)}(T)-u^{N,(m)}(T)}_H^2
  \right)^{1/2}
\]
estimates the root mean-square error. The experiment below uses
\[
  M=10000,\qquad N_{\mathrm{ref}}=4096,
  \qquad \text{random seed }271828.
\]

The theory predicts approximately
\[
  \left(\E\norm{u(T)-u^N(T)}_H^2\right)^{1/2}
  \sim N^{-1/2}
\]
for the stochastic truncation part.

\begin{table}[h]
\centering
\begin{tabular}{ccc}
\toprule
\(N\) & empirical RMS error & estimated rate \\
\midrule
\(8\) & \(0.077158\) & -- \\
\(16\) & \(0.055312\) & \(0.480\) \\
\(32\) & \(0.039370\) & \(0.490\) \\
\(64\) & \(0.027804\) & \(0.502\) \\
\(128\) & \(0.019552\) & \(0.508\) \\
\bottomrule
\end{tabular}
\caption{Monte Carlo convergence table for the spectral Galerkin approximation
with \(u_0=0\), \(T=1\), \(M=10000\) paths, and \(N_{\mathrm{ref}}=4096\).
The rate is computed between consecutive rows and is close to the predicted
root mean-square rate \(N^{-1/2}\).}
\end{table}

\subsection{Energy check}

For \(u_0=0\), the exact modal variance is
\[
  \E|u_k(t)|^2
  =
  \frac{1-e^{-2\lambda_k t}}{2\lambda_k}.
\]
Therefore
\[
  \E\norm{u^N(t)}_H^2
  =
  \sum_{k=1}^{N}
  \frac{1-e^{-2\lambda_k t}}{2\lambda_k}.
\]
This identity is a useful verification tool for simulations. The empirical
average of \(\norm{u^N(t)}_H^2\) over many paths should be close to the above
sum. In the same experiment, for \(N=128\) and \(T=1\), the empirical value was
\[
  0.082170,
\]
while the theoretical value of the finite-dimensional energy sum was
\[
  0.082939.
\]
The relative deviation is about \(0.93\%\), which is consistent with Monte
Carlo sampling error.

\subsection{Snapshot from the animation}

For the displayed paths we use the smooth polynomial initial condition
\[
  u_0(x)=10\bigl(-y^4-y^3+y^2+y\bigr),
  \qquad y=2(x-1/2).
\]
This function is \(C^\infty\), so the deterministic projection error decays
faster than the stochastic truncation error. The observed root mean-square
convergence is therefore governed by the stochastic part and follows the
\(N^{-1/2}\) rate predicted in the theorem above. In Figure
\ref{fig:animation-snapshots}, the orange curve is the deterministic heat flow
with this initial condition, while the blue curve is one stochastic
realization. The comparison illustrates that the heat semigroup damps the
deterministic profile, whereas the stochastic solution keeps receiving random
high-frequency input.

\begin{figure}[h]
\centering
\includegraphics[width=\textwidth]{spde_outputs/06_animation_snapshots.png}
\caption{Four snapshots of the deterministic heat flow and one stochastic
realization with smooth polynomial initial condition.}
\label{fig:animation-snapshots}
\end{figure}

\subsection{Expected qualitative behavior}

The numerical solution is rough in space because the driving noise is white in
space and time. Nevertheless, the heat semigroup smooths the noise enough that
the solution is \(L^2(0,1)\)-valued in one spatial dimension. Increasing \(N\)
adds more high-frequency oscillations. Their total contribution to the
mean-square \(L^2\)-error decreases like \(N^{-1}\), which corresponds to
\(N^{-1/2}\) in root mean square.

\section{Outlook}

In this project we considered the linear stochastic heat equation with additive
space-time white noise. In this special situation the spectral Galerkin
approximation leads to independent Ornstein--Uhlenbeck modes. These modes can
be simulated exactly, so that no genuine time-discretization error is
introduced. The numerical error studied here is therefore essentially the
spectral truncation error.

This is a rather favourable situation. For more general SPDEs, one usually has
to discretize time as well. For example, for semilinear equations of the form
\[
  \dd u(t) = \bigl(Au(t)+F(u(t))\bigr)\dd t + \dd W_t,
\]
the Galerkin system is no longer a collection of independent
Ornstein--Uhlenbeck processes. One therefore typically combines a spatial
discretization with time-stepping schemes such as Euler--Maruyama,
semi-implicit Euler methods, or exponential integrators. The convergence
analysis then has to control both the spatial truncation error and the
time-discretization error.

A further difficulty appears for multiplicative noise, for instance in
equations of the form
\[
  \dd u(t) = \bigl(Au(t)+F(u(t))\bigr)\dd t + G(u(t))\dd W_t .
\]
Here the stochastic integral depends on the solution itself. Thus the noise is
no longer merely an additive forcing term, but interacts directly with the
solution. If the driving noise is very rough, the solution may itself have only
distributional regularity. In such cases products like \(G(u)\xi\), where
\(\xi\) denotes space-time white noise, are not necessarily classically
defined. This is one of the central analytical difficulties in singular SPDEs.

In higher spatial dimensions this problem becomes even more pronounced, since
space-time white noise is rougher and the stochastic convolution has lower
spatial regularity. Depending on the equation, one may therefore need colored
noise, weaker solution concepts, renormalization techniques, or modern theories
such as rough paths, paracontrolled distributions, or regularity structures.

Thus the main rigorous message of this seminar is the following: in the linear
additive one-dimensional case, the stochastic heat equation can be formulated
as an abstract SPDE, solved in mild form, represented by independent
Ornstein--Uhlenbeck modes, and approximated by a spectral Galerkin method with
a transparent mean-square convergence proof. At the same time, this example
also indicates why nonlinear, multiplicative, or higher-dimensional SPDEs
require substantially more advanced analytical and numerical tools.

\appendix

\section{Ito Isometry for Cylindrical Noise}
\label{app:ito-isometry}

The following form of the Ito isometry is the basic tool used in the existence
and convergence proofs. It is standard in the Hilbert-space theory of stochastic
integration; see \cite{delSarto2025,daPratoZabczyk}.

\begin{theorem}[Ito isometry for deterministic Hilbert-Schmidt integrands]
Let \(W\) be a cylindrical Wiener process on a Hilbert space \(H\), represented
by
\[
  W(t)=\sum_{k=1}^{\infty}\beta_k(t)e_k.
\]
Let \(\Phi:[0,T]\to\mathcal L(H)\) be strongly measurable and satisfy
\[
  \int_0^T\norm{\Phi(s)}_{\HS}^2\dd s<\infty.
\]
Then
\[
  \int_0^T\Phi(s)\dd W_s
  :=
  \lim_{M\to\infty}
  \sum_{k=1}^{M}\int_0^T\Phi(s)e_k\dd\beta_k(s)
\]
exists in \(L^2(\Omega;H)\), and
\[
  \E\left\|\int_0^T\Phi(s)\dd W_s\right\|_H^2
  =
  \int_0^T\norm{\Phi(s)}_{\HS}^2\dd s.
\]
\end{theorem}

\section{Python Code for Paths and Error Computation}
\label{app:python-code}

The following code contains only the path generation and error computation used
for the numerical experiment. Plotting code is omitted.

\begin{lstlisting}
from __future__ import annotations

from dataclasses import dataclass
import numpy as np


Array = np.ndarray


@dataclass(frozen=True)
class SimulationConfig:
    n_modes: int = 96
    final_time: float = 0.25
    n_time_steps: int = 250
    seed: int | None = 1234

    @property
    def dt(self) -> float:
        return self.final_time / self.n_time_steps


def eigenvalues(n_modes: int) -> Array:
    k = np.arange(1, n_modes + 1, dtype=float)
    return (np.pi * k) ** 2


def time_grid(final_time: float, n_time_steps: int) -> Array:
    return np.linspace(0.0, final_time, n_time_steps + 1)


def ou_step_parameters(lambdas: Array, dt: float) -> tuple[Array, Array]:
    decay = np.exp(-lambdas * dt)
    variance = -np.expm1(-2.0 * lambdas * dt) / (2.0 * lambdas)
    return decay, np.sqrt(variance)


def simulate_fourier_modes(
    config: SimulationConfig,
    initial_coefficients: Array | None = None,
) -> tuple[Array, Array]:
    rng = np.random.default_rng(config.seed)
    lambdas = eigenvalues(config.n_modes)
    decay, innovation_std = ou_step_parameters(lambdas, config.dt)
    times = time_grid(config.final_time, config.n_time_steps)

    coeffs = np.empty((config.n_time_steps + 1, config.n_modes), dtype=float)
    if initial_coefficients is None:
        coeffs[0] = np.zeros(config.n_modes)
    else:
        coeffs[0] = np.asarray(initial_coefficients, dtype=float)

    for n in range(config.n_time_steps):
        noise = rng.normal(loc=0.0, scale=innovation_std, size=config.n_modes)
        coeffs[n + 1] = decay * coeffs[n] + noise

    return times, coeffs


def rms_error_table(
    n_values: Array = np.array([8, 16, 32, 64, 128]),
    n_ref: int = 4096,
    n_paths: int = 10000,
    final_time: float = 1.0,
    seed: int = 271828,
    chunk_size: int = 1000,
) -> list[tuple[int, float, float | None]]:
    rng = np.random.default_rng(seed)
    lambdas = eigenvalues(n_ref)
    variances = -np.expm1(-2.0 * lambdas * final_time) / (2.0 * lambdas)
    std = np.sqrt(variances)

    error_sq_sums = np.zeros(len(n_values))
    for start in range(0, n_paths, chunk_size):
        m = min(chunk_size, n_paths - start)
        coeffs = rng.normal(loc=0.0, scale=std, size=(m, n_ref))
        tail_sq = np.cumsum(coeffs[:, ::-1] ** 2, axis=1)[:, ::-1]
        for j, n_modes in enumerate(n_values):
            error_sq_sums[j] += np.sum(tail_sq[:, n_modes])

    rms = np.sqrt(error_sq_sums / n_paths)
    rates: list[float | None] = [None]
    for j in range(1, len(n_values)):
        rate = np.log(rms[j - 1] / rms[j]) / np.log(n_values[j] / n_values[j - 1])
        rates.append(float(rate))

    return [(int(n), float(e), r) for n, e, r in zip(n_values, rms, rates)]


if __name__ == "__main__":
    for n_modes, error, rate in rms_error_table():
        if rate is None:
            print(f"{n_modes:4d}  {error:.6f}  --")
        else:
            print(f"{n_modes:4d}  {error:.6f}  {rate:.3f}")
\end{lstlisting}

\begin{thebibliography}{9}

\bibitem{delSarto2025}
Gianmarco Del Sarto.
\emph{Stochastic Equations with Applications to Climate Sciences}.
Lecture notes, TU Darmstadt, Winter Semester 2025--2026.

\bibitem{stinnerPDE1}
Christian Stinner.
\emph{Partial Differential Equations I}.
Lecture notes, Winter term 2025/26, April 14, 2026.

\bibitem{stinnerPDE2}
Christian Stinner.
\emph{Partial Differential Equations II}.
Lecture notes, Summer term 2021, May 12, 2026.

\bibitem{daPratoZabczyk}
Giuseppe Da Prato and Jerzy Zabczyk.
\emph{Stochastic Equations in Infinite Dimensions}.
Cambridge University Press, 1992.

\bibitem{lordPowellShardlow}
Gabriel J. Lord, Catherine E. Powell, and Tony Shardlow.
\emph{An Introduction to Computational Stochastic PDEs}.
Cambridge University Press, 2014.

\bibitem{hairer2009}
Martin Hairer.
\emph{An Introduction to Stochastic PDEs}.
Lecture notes, 2009.
Available as arXiv:0907.4178 and from
\url{https://www.hairer.org/notes/SPDEs_Course.pdf}.

\end{thebibliography}

\end{document}
